\section*{Review brief: the quadratic rare-event law}
\addcontentsline{toc}{section}{Review brief}

\noindent\textbf{Purpose.} This packet asks for a correctness and prior-art
assessment of a specific statistical claim. The opening brief gives the
claim, its vulnerable dependencies, and the verification boundary. The
remaining pages contain the complete existing paper argument with its three
proof inputs expanded in place. This is an AI-assisted research draft; no
external specialist review is claimed.

\subsection*{The experiment and claimed rate}
A fresh coin has hidden magnitude $D$ with density $32d$ on $(0,1/4]$ and
hidden sign $S\mid D=d\sim\operatorname{Ber}(q(d))$. Its flips have mean
$1/2+(2S-1)D$. The unknown measurable allocation satisfies
$0\le q\le1$ and $\theta(q)=\int q\,d\mu\le\rho$, where $0<\rho\le1/4$.
Algorithms know $\mu,\rho,B$, but observe only their independent randomness
and coin flips. They may create, revisit and interleave coins, return biased
nonlinear estimates, and stop at an unbounded time. Every flip is charged.
The budget is $\sup_q\E_qT\le B$, $B\ge1$.

For $L=1+\log(1/\rho)$ and $H=1+\log L$, the minimax RMS law is
\[
 R_{\mu,\rho}(B)\asymp
 \min\left\{\rho,\sqrt{\frac{\rho LH}{B}}\right\},
\]
with universal constants for this one fixed magnitude law. The supremum is
over the full global-rarity class; it is not a pointwise constraint
$q(d)\le\rho$, and the target is not the known number $\rho$.

\subsection*{The variational conjecture and its failure}
For $h>0$, define the actual variational objective
\[
 \Psi_{\mu,\rho}(h)=
 \sup_{\substack{v\text{ measurable},\ 0\le v\le1\\\int v\,d\mu\le\rho}}
 \int\min\{v(d),h\sqrt{v(d)}/d\}\,\mu(dd).
\]
The project's earlier continuation conjectured a matching statistical upper
bound for this functional. A general lower bound is proved there using a
published single-coin input; the quadratic law refutes the matching upper.
At this law the functional has order
\[
 \Psi_{\mu,\rho}(B^{-1/2})\asymp
 \min\left\{\rho,\sqrt{\frac{\rho L}{B}}\right\}
\]
\clearpage
\section*{The exact comparison and its statistical context}

The budget $B=LH/\rho$ gives risk of order $\rho$ and objective of order
$\rho/\sqrt H$. Thus no universal matching $C\Psi$ upper bound can hold.
This is an asymptotic separation: $\sqrt H<2.54$ for
$\rho\ge10^{-100}$ in the stated rarity range, and the proof does not
establish a practically visible separation. The variational optimizer
itself remains valid.

\subsection*{Checking the variational formula}
The optimizing profile is $v_t(d)=\min\{1,t^2/d^2\}$, where $0<t<a=1/4$
is determined by its mass:
\[
 \rho=\frac{t^2}{a^2}\left(1+2\log\frac at\right).
\]
Optimality for all measurable competitors is proved by
\href{https://github.com/davidiach/adaptive-threshold-estimation/blob/eabf274cb42d4bff4daab6b58568ad6169e7f0a7/lean/ThresholdHC/Tier3/Optimizer.lean}{\texttt{arbitrary\_law\_optimizer}}.
For $h<t$, splitting at $h$ and $t$ gives
\[
 \Psi_{\mu,\rho}(h)=\frac{h^2+2h(t-h)+2ht\log(a/t)}{a^2}
 =\frac{2ht(1+\log(a/t))-h^2}{a^2}.
\]
For $h\ge t$, the objective is the feasible mass $\rho$.
These are the values in (E) below. With $u=1+2\log(a/t)$,
$\rho=ue^{1-u}$ and $L=u-\log u$; thus $u$ and $L$ are comparable.
The separate formal comparison gives constants 1 and 16 in the displayed
rate for $\Psi$.

\subsection*{Gap entropy and the contribution being assessed}
The sequential sign-testing obstruction has established antecedents:
Farrell's lower bound~\cite{farrell}, Chen--Li's sign-problem
bound~\cite{chenli}, and Chen--Li--Qiao's entropy bound~\cite[Lemma 4.1 and
Remark D.4]{clq}. The latter includes a Bernoulli extension. The finite-gap
lemma below proves the needed one-sided expected-work bound explicitly,
at fixed error and without termination under the fair coin.
Agrawal--Ram--Ramdas~\cite{agrawal} give general composite-test bounds in
the Farrell regime. Yao--Zhao--Zhou~\cite{yao} resolve the gap-entropy
conjecture for Gaussian best-arm identification and report Lean proofs.
Their observation model and target differ from this population-estimation
problem; their theorem is not used as an input here.

For intuition, a quadratic-law annulus at gap $d_j\asymp4^{-j}$ has mass
of order $16^{-j}$ and testing cost of order $16^j$. The products are
comparable across the available scales, so their normalized weights have
entropy of order $\log L$. This explains the factor $H$; it is not a new
arbitrary-law characterization. The contribution under review is the
null-cost obstruction and its transfer through one common hard pair to
the full adaptive population model, rather than a new sign-entropy principle.

The original unrestricted-sign inverse-moment boundary is at power exponent
$\alpha=1$. The rarity-dependent quadratic phenomenon here is at
$\alpha=2$. This packet treats the fixed quadratic law and does not claim
a uniform transition theorem as the exponent approaches two.

\clearpage
\section*{Proof route and questions for the reviewer}
\addcontentsline{toc}{section}{Proof route and review questions}

The lower proof uses the convex subclass
$\mathcal X_\rho=\{0\le q\le1/2,\ \int q\,d\mu\le\rho\}$ and
$J=\rho LH$. Its logical chain is as follows.

\begin{enumerate}
\item \textbf{Null-cost entropy obstruction.} For every one-coin estimator
with uniform bias at most $\rho/64$ and separately finite uniform variance
and work, the claim is
\[
 \left(\sup_{q\in\mathcal X_\rho}\Var_qY\right)\E_0T\ge c_bJ.
\]
The work is measured under the all-negative population. The proof selects
good gaps in separated annuli and restarts the estimator on unused blocks
of the \emph{same} coin to construct one sign test. Finite-prefix
change-of-measure inequalities and overlapping stopping windows supply the
entropy penalty. Fair-coin termination is not assumed.

\item \textbf{One common pair.} A finite categorical reduction and convex
RMS comparison first give a hard pair for each capped local policy.
Revealing the index of a finite policy mixture makes both its null work
and squared Hellinger distance exactly affine in the mixture weights.
Minimax and weak-star compactness are then used to obtain one separated
pair satisfying, for \emph{every} capped policy $A$,
\[
 \mathsf H^2(P^A_-,P^A_+)\le C\delta^2\E_0T_A/J.
\]
The pair is existential. It is not identified with the radial variational
optimizer or selected after observing a transcript.

\item \textbf{Population transfer and termination.} Finite tree leaves
factor by coin even under interleaving. One-coin hybrid simulations use
midpoint tapes for the other coins. The midpoint law dominates half the
null law, so the one-coin bounds can be charged to actual midpoint work.
The same endpoint pair applies at all caps. Finite expected endpoint work
then passes the testing bound to the terminal experiment.
\end{enumerate}

\noindent\textbf{Requested scrutiny.} Does the formal single-coin model
faithfully express the intended \emph{null-cost} theorem? Are the
sufficiency, convexity and compactness hypotheses satisfied strongly enough
to obtain a pair common to all policies? Does the hybrid simulation charge
work under the correct law, including arbitrary randomization and creation?
Does the endpoint limit preserve the full unrestricted model?

Separately, is an equivalent critical population-estimation result already
known? The complete proof credits Chen--Li--Qiao, Juditsky--Nemirovski,
Polyanskiy--Wu and Lee--Valiant. Their ingredients and this particular
application should be assessed separately.

\clearpage
\begingroup\small
\section*{What has been checked, and what remains open}
\addcontentsline{toc}{section}{Verification and provenance}

\begin{center}
\begin{tabular}{>{\raggedright\arraybackslash}p{0.34\textwidth}>{\raggedright\arraybackslash}p{0.58\textwidth}}
\toprule
Claim or component & Verification boundary\\
\midrule
Original arbitrary-law RMS law & Both directions are Lean-checked,
with coefficients $1/64$ and $512$. One budget-indexed upper family
works for every gap law in the original observable experiment.\\[5pt]
Single-coin base cost & The BC1 variance/null-work bound and its raw-moment
version are Lean-checked for all rarities. The proof derives activation,
good gaps, observable restarts and the BC2 entropy bound from actual
seeded Bernoulli experiments.\\[5pt]
Critical minimax lower bound & Lean-checked for the full rare-prior class.
The proof constructs one common pair, transfers it to unrestricted
population algorithms, and proves endpoint termination and testing.
Budget nonemptiness and a trivial upper bound are also proved.\\[5pt]
Critical upper and $\Psi$ comparison &
\href{https://github.com/davidiach/adaptive-threshold-estimation/blob/eabf274cb42d4bff4daab6b58568ad6169e7f0a7/lean/ThresholdHC/CriticalUpper/Full.lean}{Lean-checked upper}
through dyadic classifiers, coupled corrections on one prefix and fixed-count
aggregation, with coefficient $100000$ on the \emph{RMS} scale. The written
appendix instead uses a likelihood bank and Poisson replication. Also checked:
\href{https://github.com/davidiach/adaptive-threshold-estimation/blob/eabf274cb42d4bff4daab6b58568ad6169e7f0a7/lean/ThresholdHC/CriticalPsi/Full.lean}{exact $\Psi$ formula, rate and counterexample}.
The formal witness uses $B=LH/\rho$; the paper's $\rho^{-2}$ sequence
is a separate derivation. No external specialist review is claimed.\\[5pt]
Numerical and exact finite checks & Sanity and reproducibility evidence.
They do not certify the continuum theorem or distinguish its asymptotic
rate from the proposed Psi rate.\\
\bottomrule
\end{tabular}
\end{center}

The existing critical proof was reviewed by an adversarial agent. A second
review disclosed that its author also wrote the base-cost component; it
must not be counted as a fully independent review of that component.
The recorded construction-era reviews do not establish model diversity. The follow-up audit has not identified
a blocking error in its initial pass, but that observation is not a
certificate of correctness.

\subsection*{Source and reproduction}
The complete mathematical argument below is expanded from the four
source files at published commit
\texttt{b3f88e6bc02b3d59ec9fe99529688e2bd911e1db}.
The accompanying build script verifies their existing SHA-256 lock before
assembly. Editorial changes concern the title, this brief, additional references,
the historical status paragraph, and layout. Inherited mathematical statements and arguments are preserved;
monospaced source paths and hashes are made line-breakable for typesetting.
The verification table refers to the later formal checkpoint
\nolinkurl{eabf274cb42d4bff4daab6b58568ad6169e7f0a7}.

The checked toolchain is Lean \texttt{v4.36.0-rc1}, with Mathlib revision
\nolinkurl{9bdd0bc71aaae0fd83889f9aacd70c32fccd1bcd}.
The audits report only \texttt{propext}, \texttt{Classical.choice} and
\texttt{Quot.sound}; their exact target types and axiom output are linked
from the verification guide. The observable model includes the independent
seed and a finite flip history, allows arbitrary natural-number coin labels,
and assigns infinite work to nontermination. Finite expected work gives
almost-sure termination separately under each endpoint law.

\begin{itemize}
\item \href{https://github.com/davidiach/adaptive-threshold-estimation/blob/eabf274cb42d4bff4daab6b58568ad6169e7f0a7/docs/CRITICAL_AUDIT.md}{Follow-up dependency audit and formalization boundary}.
\item \href{https://github.com/davidiach/adaptive-threshold-estimation/tree/b3f88e6bc02b3d59ec9fe99529688e2bd911e1db/research/2026-10-27-nonlinear-critical}{Original proof bundle, source lock, and scoped reviews}.
\item \href{https://github.com/davidiach/adaptive-threshold-estimation/blob/6009662a5581b31de17c8956b6fd93e4279ba087/docs/process/METHOD.md}{Process note and limits of the evidence}.
\end{itemize}

\noindent A matching rare-event characterization for arbitrary magnitude laws remains
open. This packet concerns a fixed quadratic law. It does not establish
uniformity as a power exponent approaches two, explicit useful constants,
or practical superiority of an implementation.

\endgroup
\clearpage
\part*{Complete critical-law argument}
\addcontentsline{toc}{part}{Complete critical-law argument}
